
\section{Les conditions en C}

En C, les conditions permettent d'exécuter certaines parties du programme uniquement si des critères sont vérifiés.

\subsection{Structure conditionnelle if / else}
\begin{UPSTIinfor}{Structure conditionnelle}
	Une structure conditionnelle permet d'exécuter une partie du programme uniquement si une condition est vérifiée. Dans le cas contraire, une autre partie du programme peut être exécutée ou le programme peut continuer son exécution normalement. 
	On peut avoir une succesion de conditions, chacune étant testée à son tour si la précédente n'est pas vérifiée. Le programme cesse d'évaluer les conditions dès qu'une condition est vérifiée.

	Une structure conditionnelle est composée de :
	\begin{itemize}
		\item une condition (une expression qui peut être vraie ou fausse)
		\item un bloc de code à exécuter si la condition est vraie
		\item éventuellement d'autres conditions à tester si la condition précédente est fausse
		\item éventuellement un bloc de code à exécuter si toutes les conditions sont fausses
	\end{itemize}
	
\end{UPSTIinfor}

En langage C, la structure conditionnelle est écrite de la manière suivante :

\begin{lstlisting}[language=c]
if (<condition1>) // Si la condition 1 est vraie
{
    // Partie exécutée si la <condition1> est vérifiée
}
else if (<condition2>) // Sinon, si la condition 2 est vraie
{
    // Partie exécutée si la <condition2> est vérifiée
}
else // Sinon, dans tous les autres cas
{
    // Partie exécutée par défaut
}
\end{lstlisting}

Exemple d'utilisation de la structure conditionnelle
	\begin{lstlisting}[language=c]
int x = 5;
if (x > 0)
{
	printf("x est positif\n");	
}
else if (x < 0)
{
	printf("x est négatif\n");
}
else
{
	printf("x vaut zéro\n");
}
	\end{lstlisting}

La structure \texttt{if} peut prendre plusieurs formes :

\begin{itemize}
	\item \texttt{if} seul
	\item \texttt{if} puis \texttt{else}
	\item \texttt{if} puis \texttt{else if} puis \texttt{else}
\end{itemize}

On peut placer autant de \texttt{else if} que l'on veut à la suite.
\subsection{Opérateurs de comparaison}
\label{sec:operateur-comparaison}

Les comparaisons sont faites à l'aide des opérateurs suivants :

\begin{center}
	\begin{tabular}{|l|l|l|l|}
		\hline
		Symbole     & Signification     & Exemple (\texttt{x=5}, \texttt{y=10}) & Résultat \\
		\hline
		\texttt{==} & égal              & \texttt{x == 5}                       & vrai     \\
		\texttt{!=} & différent         & \texttt{x != y}                       & vrai     \\
		\texttt{<}  & inférieur         & \texttt{x < y}                        & vrai     \\
		\texttt{>}  & supérieur         & \texttt{x > y}                        & faux     \\
		\texttt{<=} & inférieur ou égal & \texttt{x <= 5}                       & vrai     \\
		\texttt{>=} & supérieur ou égal & \texttt{y >= 11}                      & faux     \\
		\hline
	\end{tabular}
\end{center}
\begin{UPSTIwarning}{Piège ! Ne pas confondre affectation et comparaison}
	L'opérateur \texttt{=} sert à affecter une variable et l'opérateur \texttt{==} sert à comparer deux variables ! Si vous les confondez, rien ne fonctionnera...
\end{UPSTIwarning}

\subsection{Opérateurs logiques}

On peut combiner plusieurs conditions avec les opérateurs logiques :

\begin{center}
	\begin{tabular}{|l|l|l|l|}
		\hline
		Symbole                                         & Signification & Exemple (x=5, y=10)                                            & Résultat attendu \\
		\hline
		\texttt{\&\&} & ET logique & \texttt{(x > 0 \&\& y < 20)}  & vrai \\
		\texttt{||}   & OU logique & \texttt{(x < 0 || y > 100)}   & vrai \\
		\texttt{!}    & NON logique & \texttt{!(x == 5)}             & faux \\
		\hline
	\end{tabular}
\end{center}
\begin{lstlisting}[language=c]
if ((x > 0) && (x < 10))
{
    printf("x est compris entre 0 et 10\n");
}

if ((x < 0) || (x > 100))
{
    printf("x est en dehors de l'intervalle [0,100]\n");
}
\end{lstlisting}

